% Bilinear rank as proof cost: fast matmul schemes and field-specific
% flip search for zkML.  Build: tectonic doc/matmul_zkml_paper.tex
\documentclass[10pt]{article}
\usepackage[margin=2.4cm]{geometry}
\usepackage{amsmath,amssymb}
\usepackage{mathtools}
\usepackage{booktabs}
\usepackage{graphicx}
\usepackage{fancyvrb}
\usepackage{enumitem}
\usepackage{microtype}
\usepackage[hidelinks]{hyperref}
\urlstyle{same}
\setlist{itemsep=2pt,topsep=3pt}
\setlength{\emergencystretch}{3em}
\DefineVerbatimEnvironment{Code}{Verbatim}{fontsize=\footnotesize}

\title{\bfseries Bilinear Rank as Proof Cost: Fast
  Matrix-Multiplication Schemes and Field-Specific Flip Search for
  zkML}
\author{Greg Sidebottom}
\date{}

\begin{document}
\maketitle
\vspace{-3em}
\begin{center}
\itshape Research note --- logic repo
(\url{https://github.com/gsidebottom/logic}), matmul track,
2026-07-11. Companion to the $3\times3$ additive-complexity paper
(\texttt{doc/matmul\_adds\_paper.md}, artifacts DOI
10.5281/zenodo.21240904), the rank-48 rigidity paper, and the flower
note. The engine of \S\ref{sec:engine} is
\texttt{src/bin/flip23p.rs} in this repository.
\end{center}

\begin{abstract}
Zero-knowledge machine learning (zkML) --- proving that a
neural-network inference was computed correctly without re-executing
or revealing it --- is bottlenecked by the prover, whose work in
circuit-based proof systems is counted in \textbf{multiplication
constraints}: additions and constant-coefficient linear combinations
are free. This is the exact inversion of the hardware cost model that
keeps fast matrix multiplication out of GPUs, and it makes
\textbf{bilinear rank} --- the number of multiplications in a
Strassen-like scheme --- the direct cost driver for matrix products
between witness values (private weights; attention). Three
consequences. (1) Known schemes already pay: recursive rank-48
$4\times4$ blocking proves a $768\times768$ witness$\times$witness
product with ${\approx}4\times$ fewer constraints than the naive
circuit, penalty-free because the schemes' additions cost nothing.
(2) Coefficient obstructions vanish: schemes over $\mathbb{Q}$ with
dyadic (or any rational) coefficients are exact over a prime field,
and numerical stability is not a concept. (3) Most interestingly,
\textbf{rank is field-dependent} --- rank 47 for $4\times4$ exists
modulo 2 with no characteristic-0 equivalent known --- and zkML fixes
a short list of concrete fields (Goldilocks, BabyBear, M31, the
BN254/BLS12-381 scalar fields) over which nobody has searched. We
port our rational flip-graph engine to $\mathbb{F}_p$
(\texttt{flip23p}, Goldilocks first): over a prime field the engine
is strictly stronger than its $\mathbb{Q}$ parent --- no coefficient
growth, no magnitude caps, and every coincidence-solving $\lambda$ is
admissible, so the coplanarity structure that rationality left mostly
unusable becomes fully spendable. A verified rank-22 $3\times3$ or
rank-47 $4\times4$ over a proof field would be a shippable
constraint-count reduction with no analog in the characteristic-0
literature. We state the cost model precisely, including the honest
carve-outs (sumcheck/GKR provers multiply matrices in ${\sim}n^2$
prover work and do not care about rank; public-weight linear layers
cost no constraints), inventory which of our existing artifacts
transfer, and report the port's first measurements.
\end{abstract}

\section{zkML in one page}\label{sec:zkml}

A \textbf{zero-knowledge proof} (ZKP) lets a prover convince a
verifier that a statement holds --- here, ``this output is what model
$M$ computes on input $x$'' --- without the verifier re-executing the
computation, and optionally without revealing $x$ or $M$. Two
families dominate. \textbf{SNARKs} (succinct non-interactive
arguments of knowledge; Groth16 \cite{groth16}, PLONK \cite{plonk})
give constant-or-polylog proof sizes over elliptic-curve-pairing
fields such as the BN254 and BLS12-381 scalar fields (${\sim}254$-bit
primes). \textbf{STARKs} \cite{stark} replace pairings with hashes,
gain transparency and plausible post-quantum security, and run over
small ``FFT-friendly'' primes chosen for fast number-theoretic
transforms --- Goldilocks $p = 2^{64} - 2^{32} + 1$ (Plonky2
\cite{plonky2}), BabyBear $p = 2^{31} - 2^{27} + 1$ (RISC~Zero
\cite{risc0}), and M31 $= 2^{31}-1$ (Circle STARKs \cite{circle}).

The asymmetry that defines the field: verification is cheap,
\textbf{proving is $10^3$--$10^6\times$ the native computation}, all
of it exact arithmetic in the proof field. zkML applications ---
proving a proprietary model was faithfully executed (model privacy),
proving an inference inside a blockchain rollup, auditable AI where
the checker is cheap --- inherit that overhead on every multiply of
every layer. Making matrix multiplication cheap \emph{inside a proof}
is therefore a different optimization problem from making it fast on
silicon, with a different --- in fact opposite --- cost model.

\section{What proving charges}\label{sec:cost}

The standard arithmetization, \textbf{R1CS} (rank-1 constraint
systems, as in Groth16 \cite{groth16}), expresses a computation as
constraints of the form
\[
\Big( \textstyle\sum_i a_i w_i \Big) \times
\Big( \textstyle\sum_i b_i w_i \Big) =
\Big( \textstyle\sum_i c_i w_i \Big)
\]
over the proof field, where $w$ is the witness vector and the $a_i,
b_i, c_i$ are \emph{constants}. Each constraint carries exactly one
multiplication of witness values; the linear combinations inside the
parentheses --- additions, subtractions, scalings by arbitrary field
constants --- are free: they are folded into the constraint's
constant vectors and cost the prover nothing. PLONK-family systems
\cite{plonk} charge per gate/row with the same shape (custom gates
and lookup arguments shift constants around but preserve the
headline: \textbf{bilinear multiplications of witnesses are the
metered resource}).

For an $n\times n$ product $C = A\cdot B$ where both operands are
witnesses, each entry $C_{ij} = \sum_k A_{ik}B_{kj}$ is a sum of $n$
witness$\times$witness products; a sum of $n$ products is not one
rank-1 quadratic form --- it takes $n$ constraints. The naive circuit
therefore costs $n^3$ constraints. Two situations make both operands
witnesses in zkML: \textbf{private weights} (the commercially central
case --- the model owner proves inference without revealing the
model) and \textbf{attention} ($QK^{\mathsf T}$ and
$\mathrm{scores}\cdot V$ multiply two activation matrices ---
witness$\times$witness even when all weights are public).
Conversely, a linear layer with \emph{public} weights is a constant
linear map and costs zero multiplication constraints; rank
optimization has nothing to offer there.

\paragraph{A toy example, $2\times2$ over $\mathbb{F}_{17}$.} Take
witnesses $A = \begin{psmallmatrix}1&2\\3&4\end{psmallmatrix}$,
$B = \begin{psmallmatrix}5&6\\7&0\end{psmallmatrix}$; then
$C = A\cdot B \equiv \begin{psmallmatrix}2&6\\9&1\end{psmallmatrix}
\pmod{17}$. The naive circuit spends one constraint per product ---
eight constraints of the shape $u_1 = A_{11}\times B_{11}$ --- and
output sums like $C_{11} = u_1 + u_2$ cost nothing further (they are
linear). Strassen's rank-7 scheme spends \textbf{seven}: its first
product is the single constraint
\[
( w_{A_{11}} + w_{A_{22}} ) \times ( w_{B_{11}} + w_{B_{22}} )
  = w_{P_1},
\]
where all four operand additions sit \emph{inside} the constraint's
linear forms, free. There is exactly \textbf{one witness per proof}:
a single vector $w$ over $\mathbb{F}_{17}$ holding every wire value
--- a conventional leading 1 (which lets constraints use constants),
the inputs, the seven products, the outputs. Concretely, all 20
coordinates:
\[
w = (\,1 \mid \underbrace{1,2,3,4}_{A} \mid \underbrace{5,6,7,0}_{B}
      \mid \underbrace{8,1,6,8,0,5,3}_{P_1..P_7}
      \mid \underbrace{2,6,9,1}_{C}\,).
\]
The subscripted symbols above are coordinates of this one vector,
and each constraint is checked by reading them off: for $P_1$,
$(1+4)\cdot(5+0) = 25 \equiv 8$ \checkmark, and likewise
$P_2 = 7\cdot5 = 1$, $P_3 = 1\cdot6 = 6$, $P_4 = 4\cdot2 = 8$,
$P_5 = 3\cdot0 = 0$, $P_6 = 2\cdot11 = 5$, $P_7 = (-2)\cdot7 = 3$
(all mod 17). The free output combinations reproduce $C$ exactly:
$C_{11} = P_1{+}P_4{-}P_5{+}P_7 = 19 \equiv 2$,
$C_{12} = P_3{+}P_5 = 6$, $C_{21} = P_2{+}P_4 = 9$,
$C_{22} = P_1{-}P_2{+}P_3{+}P_6 = 18 \equiv 1$. (The prover commits
to $w$ and proves the constraints hold; the verifier never sees $w$
--- that is the ``zero-knowledge''. Note the naive circuit's witness
carries eight product coordinates to Strassen's seven, so lower rank
also shortens the committed vector.) Seven constraints instead of
eight --- 12.5\% at one level, compounding to $n^{2.807}$ under
recursion --- and dyadic coefficients would be equally at home:
$\tfrac12$ is just the constant 9 in $\mathbb{F}_{17}$
($2\cdot9 = 18 \equiv 1$). This is \S\ref{sec:rank}'s story in
miniature; the rank-23 and rank-48 schemes are the same picture with
better ratios. (Appendix~\ref{app:tutorial} runs a complete Setup
$\to$ Prove $\to$ Verify on a two-constraint slice of this example
--- the whole SNARK pipeline at hand-checkable scale.)

\section{Bilinear rank is constraint count}\label{sec:rank}

A bilinear scheme of rank $r$ computes an $m\times m$ block product
with $r$ products of linear forms. In R1CS each product is exactly
one constraint --- the linear forms slide into the constraint's free
linear parts --- so an $m\times m$ witness$\times$witness block
product costs $r$ constraints instead of $m^3$, and recursion
compounds it: $T(n) = r\,T(n/m)$, i.e.\ $n^{\log_m r}$ constraints
total, \textbf{with zero cost for the combination additions at every
level}. The additions that price fast schemes out of silicon are
literally free here; there is no numerical-stability tax because
arithmetic is exact; and coefficients may be any field constants ---
dyadic $\tfrac12$'s are simply the constant $(p+1)/2 \bmod p$.

Idealized constraint counts for a witness$\times$witness product
(recursion to scalar base; real circuits stop at a base tile, which
shifts absolute numbers but not the ordering):

\begin{center}
\begin{tabular}{lcccc}
\toprule
scheme (exponent) & $n=768$ & vs naive & $n=4096$ & vs naive\\
\midrule
naive (3.000) & $4.53\times10^{8}$ & $1.0\times$ & $6.87\times10^{10}$ & $1.0\times$\\
$3{\times}3:23$ (2.854) & $1.72\times10^{8}$ & $2.6\times$ & $2.05\times10^{10}$ & $3.4\times$\\
Strassen $2{\times}2:7$ (2.807) & $1.26\times10^{8}$ & $3.6\times$ & $1.39\times10^{10}$ & $4.9\times$\\
\textbf{$4{\times}4:48$ (2.792)} & $1.14\times10^{8}$ & $\mathbf{4.0\times}$ & $1.23\times10^{10}$ & $\mathbf{5.6\times}$\\
$3{\times}3:22$ (2.814), \emph{if found} & $1.32\times10^{8}$ & $3.4\times$ & --- & ---\\
$4{\times}4:47$ (2.785), \emph{if found} & $1.09\times10^{8}$ & $4.2\times$ & --- & ---\\
\bottomrule
\end{tabular}
\end{center}

Two readings. First, \textbf{the wins in the top half are available
today}: the rank-48 scheme \cite{alphatensor,dps} with dyadic
coefficients drops into any R1CS/PLONK matmul gadget as-is (our
repository carries the verified scheme and its minimized linear
networks). Second, the \emph{marginal} rows show why new records
matter here more than in silicon: each rank step is a few percent per
recursion level, compounding --- and unlike hardware, nothing eats
the margin.

\section{The honest carve-outs}\label{sec:carve}

Rank is \textbf{not} the lever everywhere:

\begin{itemize}
\item \textbf{Sumcheck/GKR provers do matmul in ${\sim}n^2$ prover
  field-ops} (Thaler \cite{thaler}; GKR \cite{gkr}); zkCNN
  \cite{zkcnn} and successors build dedicated zkML provers on this
  line, and where such a prover applies, bilinear rank is irrelevant.
  The rank lever applies to the (large) circuit-based world:
  general-purpose SNARK toolchains, zkVMs executing matmul as
  arithmetic, PLONK/R1CS pipelines such as EZKL \cite{ezkl}, and
  folding-based IVC (Nova \cite{nova}) whose per-step cost is the
  step circuit's constraint count.
\item \textbf{Public-weight linear layers are free}; rank helps
  witness$\times$witness products only --- private weights and
  attention.
\item \textbf{Constraint count is not wall-clock}: prover time also
  includes commitments (multi-scalar multiplications / hashing);
  fewer constraints shrink those too, but constants differ by
  system.
\item \textbf{Nonlinearities dominate some workloads} (range checks,
  lookups for ReLU/quantization); the matmul share is large in
  attention/convolution-heavy models, smaller in lookup-heavy
  quantized pipelines.
\end{itemize}

\section{Field-specific rank: the open frontier, and
  flip23p}\label{sec:engine}

Bilinear rank depends on the coefficient field. The emblem:
AlphaTensor's rank-\textbf{47} $4\times4$ scheme exists over
$\mathbb{F}_2$ \cite{alphatensor}, no characteristic-0 rank-47 is
known, and our companion work proved the only $\mathbb{Q}$-usable
rank-48 class \emph{rigid} under a large certified move system ---
evidence that characteristic 0 is genuinely constrained in ways a
fixed finite field need not be. zkML nominates a short, concrete
list of fields that matter commercially --- Goldilocks, BabyBear,
M31, BN254-Fr, BLS12-381-Fr --- and, to our knowledge, \textbf{no
rank search has ever been run over any of them}: the flip-graph
literature works over $\mathbb{F}_2$ and small extensions
\cite{km}, the explicit records over $\mathbb{Z}/\mathbb{Q}$. A
verified $3\times3$ rank-22 --- or $4\times4$ rank-47 --- over
Goldilocks would be a new kind of record with an immediate consumer:
2--4\% fewer constraints per recursion level in deployed proof
systems, compounding as in \S\ref{sec:rank}.

\texttt{flip23p} is our rational flip engine (splits,
$\lambda$-flips, reductions, canonical gauge, exact Brent
verification --- the machinery of the companion papers) re-based
onto $\mathbb{F}_p$, Goldilocks first. Over a prime field the engine
is \emph{strictly stronger} than its $\mathbb{Q}$ parent:
\begin{enumerate}
\item \textbf{No coefficient growth.} $\mathbb{Q}$-walks need
  magnitude caps and dyadic bookkeeping; field elements don't grow.
  Every cap in the $\mathbb{Q}$ engine is deleted, not loosened.
\item \textbf{Every solved flip is admissible.} Over $\mathbb{Q}$, a
  coincidence ($f_i + \lambda f_j \propto f_m$, the Cramer-solved
  move that makes flips productive) is usable only when $\lambda$ is
  dyadic --- the filter discarded most solutions. Over $\mathbb{F}_p$
  every nondegenerate solution is a legal move: the coplanarity
  structure (``H4's raw material'' in the flower note) becomes fully
  spendable.
\item \textbf{The gauge simplifies}: monic normalization (leading
  coefficient 1, scalars folded into the $c$-slot); proportional
  $\Leftrightarrow$ equal, with none of $\mathbb{Q}$'s content/sign
  bookkeeping.
\end{enumerate}

First measurements (2026-07-11; \S\ref{sec:repro}): the 55-addition
record scheme loads and verifies over Goldilocks (729 Brent equations
mod $p$); the census gate matches the $\mathbb{Q}$ engine exactly (9
shared pairs, weight 177, distinct-rows 40). One honest surprise: the
seed's solved-move silence is \emph{not} the dyadic filter's fault
--- even with all field $\lambda$ available, its 9 shared pairs admit
no coincidence solution; the coplanarity geometry genuinely fails
there. Mobility arrives at once elsewhere: a 25-second smoke storm
from the most mobile census seed executed 1.3M moves with 369K solved
flips and ${\sim}$10K reductions --- a far higher solved-flip
fraction than the $\mathbb{Q}$ storm ever achieved, the freed
$\lambda$ paying immediately. No sub-23 rank has appeared in smoke
tests (nor was one expected at that effort); the campaigns are the
point.

\section{What transfers from the existing program}\label{sec:assets}

\begin{itemize}
\item \textbf{Engines and discipline.} flip48/flip23's storm,
  closure, census, novelty and thin-storm machinery ports
  mechanically (v1 of \texttt{flip23p} carries
  storm/census/closures; the beam-chase and exhaustive campaign
  modes follow). The operational rules hard-won in the companion
  papers --- completion logs, full-frontier dumps, control
  baselines, budget caps --- apply verbatim.
\item \textbf{Seeds.} Every scheme we hold is a valid $\mathbb{F}_p$
  scheme: the 17,376 database classes, our 53 lifted classes, the 46
  storm-new classes, and the 467 mod-2-invisible dyadic
  $\mathbb{Q}$-schemes (denominators are field constants). The
  mobility census that ranked $\mathbb{Q}$-seeds re-runs unchanged.
\item \textbf{Verified gadgets.} The Lean pipeline that certified
  the 55-add scheme extends naturally: a scheme-to-circuit emitter
  (circom / Halo2 templates for the 23- and 48-multiplication
  blocks) with a machine-checked correctness certificate is a
  deliverable the proof-systems world actually values.
\item \textbf{What does not transfer:} the additive-complexity
  results (55 adds, the no-54 theorem). Additions are free here;
  that program's value stays in classical hardware and in
  witness-generation speed.
\end{itemize}

\section{Program}\label{sec:program}

(1) Goldilocks campaigns: storm + closure portfolios over the
census-ranked seeds; targets, in order of plausibility $\times$
value: any new-class harvest over $\mathbb{F}_p$, rank 22 at
$3\times3$, rank 47 at $4\times4$ (the $4\times4$ port is the same
edit at 16 coordinates). (2) Small-field variants: BabyBear and M31
(31-bit arithmetic, ${\sim}4\times$ faster walks), then BN254-Fr.
(3) Rectangular tiles: transformer shapes want $\langle
m,n,k\rangle$ records; the engine generalizes as flip48 $\to$ flip23
did. (4) Verified gadget emitter + benchmark note: constraint-count
measurements in an EZKL-style stack, naive vs recursive-scheme, with
Lean-certified gadget templates.

\section{Reproduction}\label{sec:repro}

\begin{Code}
cargo build --release --bin flip23p
./target/release/flip23p --census --dir matmul/mm23      # gates
./target/release/flip23p --dir matmul/seeds23/13a5bd4n \
    --seconds 25 --threads 4 --out matmul/found23p       # smoke storm
# any verified rank <= 22 over Goldilocks lands in
# found23p/RECORDP_rank*.txt and is announced loudly
\end{Code}

\appendix
\section{Setup, Prove, Verify --- a complete toy proof}\label{app:tutorial}

This appendix runs the entire SNARK pipeline, Groth16-shaped, on a
two-constraint slice of \S\ref{sec:cost}'s example, every number over
$\mathbb{F}_{17}$. One honest simplification, flagged where it
occurs: real systems wrap the values below in elliptic-curve
encodings (``in the exponent''), which is what makes hiding
cryptographic; the \emph{arithmetic} --- R1CS $\to$ QAP $\to$
divisibility check at a secret point --- is shown exactly as real
provers compute it.

\paragraph{The statement.} For reference, the full toy matrices of
\S\ref{sec:cost}, whose entries supply every private value below:
\[
A = \begin{pmatrix} 1 & 2\\ 3 & 4 \end{pmatrix}, \qquad
B = \begin{pmatrix} 5 & 6\\ 7 & 0 \end{pmatrix}.
\]
The prover claims: ``I know private values $A_{11}, A_{22}, B_{11},
B_{12}, B_{22}$ (here $1, 4, 5, 6, 0$ --- the marked entries of $A$
and $B$) such that $P_1 = (A_{11}{+}A_{22})(B_{11}{+}B_{22}) = 8$ and
$P_3 = A_{11}(B_{12}{-}B_{22}) = 6$.'' The claimed products $(8,6)$
are \textbf{public}; the five inputs stay private. The witness
(public part first):
\[
w = (\,1,\ P_1{=}8,\ P_3{=}6 \mid A_{11}{=}1,\ A_{22}{=}4,\
      B_{11}{=}5,\ B_{12}{=}6,\ B_{22}{=}0\,)
\]
and the two constraints (check: $5\cdot5 = 25 \equiv 8$,
$1\cdot6 = 6$):
\begin{align*}
c_1&: (w_{A_{11}} + w_{A_{22}}) \times (w_{B_{11}} + w_{B_{22}})
      = w_{P_1}\\
c_2&: (w_{A_{11}}) \times (w_{B_{12}} - w_{B_{22}}) = w_{P_3}
\end{align*}

\paragraph{Step 0 --- constraints become one polynomial identity
(the QAP).} Assign constraint $c_j$ to the point $x = j$, and
interpolate each coordinate's coefficients across constraints with
the basis $L_1(x) = 2-x$, $L_2(x) = x-1$:

\begin{center}
\begin{tabular}{lcccccc}
\toprule
coord & A-side of & A-poly & B-side of & B-poly & C-side of & C-poly\\
\midrule
$A_{11}$ & $c_1,c_2$ & $1$ & --- & 0 & --- & 0\\
$A_{22}$ & $c_1$ & $2-x$ & --- & 0 & --- & 0\\
$B_{11}$ & --- & 0 & $c_1$ & $2-x$ & --- & 0\\
$B_{12}$ & --- & 0 & $c_2$ & $x-1$ & --- & 0\\
$B_{22}$ & --- & 0 & $c_1 (+),\, c_2 (-)$ & $3-2x$ & --- & 0\\
$P_1$ & --- & 0 & --- & 0 & $c_1$ & $2-x$\\
$P_3$ & --- & 0 & --- & 0 & $c_2$ & $x-1$\\
\bottomrule
\end{tabular}
\end{center}

Weight each polynomial by its witness value and sum:
\[
A(x) = 1\cdot1 + 4(2-x) = 9 - 4x, \qquad
B(x) = 5(2-x) + 6(x-1) = x + 4, \qquad
C(x) = 8(2-x) + 6(x-1) = 10 - 2x.
\]
Sanity: at $x{=}1$: $A\!\cdot\!B = 25 \equiv 8 = C$ ($c_1$
\checkmark); at $x{=}2$: $A\!\cdot\!B = 6 = C$ ($c_2$ \checkmark).
\textbf{Both constraints hold $\iff$ $A(x)B(x) - C(x)$ vanishes at
$x = 1, 2$ $\iff$ it is divisible by $Z(x) = (x-1)(x-2)$.} Over
$\mathbb{F}_{17}$:
\[
P(x) = A B - C = -4x^2 - 5x + 26 \equiv 13x^2 + 12x + 9,
\qquad Z(x) = x^2 - 3x + 2,
\qquad H(x) = P/Z = 13
\]
(exact division, remainder 0). The secret knowledge is compressed
into: ``I can exhibit witness-weighted $A, B, C$ and a quotient $H$
with $AB - C = HZ$.''

\paragraph{Step 1 --- Setup (the trusted ceremony).} Sample a secret
point, say $\boldsymbol{\tau = 5}$ --- the ``toxic waste'' --- and
publish the powers of $\tau$ and the QAP polynomials evaluated at
$\tau$, each wrapped in an encoding that supports additions and
scalings but hides the value (elliptic-curve points $g^{\tau^i}$ in
real systems; bare numbers here, with hiding flagged). Then
\textbf{destroy $\tau$}: anyone holding it can forge proofs ---
hence the multi-party ceremonies around Groth16, and the transparent
(no-$\tau$) alternatives like STARKs. Published for the toy:
$Z(5) = 4\cdot3 = 12$ and the encoded coordinate-polynomial
evaluations at 5.

\paragraph{Step 2 --- Prove.} The prover, holding $w$, computes the
\emph{encoded} evaluations --- each a \textbf{linear combination} of
published encodings with witness coefficients (this is why R1CS's
free-linear-combination structure is the whole design):
\[
A(5) = 9 - 20 \equiv 6, \qquad B(5) = 9, \qquad C(5) = 0,
\qquad H(5) = 13.
\]
The proof is $\pi = (\mathrm{enc}(A(5)), \mathrm{enc}(B(5)),
\mathrm{enc}(C_{\text{priv}}(5)), \mathrm{enc}(H(5)))$ --- in
Groth16, after optimizations, three curve points, ${\sim}200$ bytes,
\emph{independent of circuit size}. The prover never learns $\tau$.

\paragraph{Step 3 --- Verify.} The verifier reconstructs the
\emph{public} part of $C(\tau)$ itself from the claimed outputs
$(8, 6)$ --- the prover cannot lie about what is being claimed ---
then checks \textbf{one multiplicative relation} (one pairing
equation in real systems; pairings are the tool that multiplies two
hidden values exactly once):
\[
A(\tau)B(\tau) - C(\tau) \stackrel{?}{=} H(\tau)Z(\tau):
\qquad 6\cdot9 - 0 = 54 \equiv 3, \qquad 13\cdot12 = 156 \equiv 3.
\quad\checkmark
\]

\paragraph{Why this convinces (soundness).} A cheating prover
without a valid witness needs $AB - C$ divisible by $Z$; the best it
can do is fake the relation at the single hidden point $\tau$. Two
distinct low-degree polynomials agree at a random point with
probability $\le \deg/|\mathbb{F}|$ --- about $2/17$ in the toy
(which is why $\mathbb{F}_{17}$ is a classroom field), about
$2^{-250}$ over a real 254-bit field. Not knowing $\tau$, the prover
cannot aim: Schwartz--Zippel is the heart of the construction.

\paragraph{Why it reveals nothing (zero-knowledge).} The verifier
sees only encodings, never witness coordinates. (Full Groth16
additionally randomizes each proof with masking multiples of $Z(x)$
so two proofs of the same witness look independent; omitted for
clarity.)

\paragraph{The connection to rank, one last time.} Each constraint
became one interpolation point, one row of the QAP table, one share
of prover work and setup size. Prove a witness$\times$witness matrix
product with a rank-$r$ scheme and this whole pipeline is $r$
constraints per block instead of $m^3$ --- the scheme's additions
never appear as constraints at any stage; they live inside the
linear combinations, which the encodings support for free.

\section*{Acknowledgments}

This work was carried out in an extended interactive collaboration
with Claude (Anthropic; the Fable 5 and Opus 4.8 models), which
implemented the engines and drafted this text under the author's
direction and review.

\begin{thebibliography}{19}
\bibitem{strassen} V. Strassen. \emph{Gaussian elimination is not
  optimal.} Numer. Math. 13:354--356, 1969.
\bibitem{stark} E. Ben-Sasson, I. Bentov, Y. Horesh, M. Riabzev.
  \emph{Scalable, transparent, and post-quantum secure computational
  integrity.} ePrint 2018/046.
\bibitem{alphatensor} A. Fawzi et al. \emph{Discovering faster matrix
  multiplication algorithms with reinforcement learning.} Nature
  610:47--53, 2022.
\bibitem{thaler} J. Thaler. \emph{Time-optimal interactive proofs for
  circuit evaluation.} CRYPTO 2013.
\bibitem{plonk} A. Gabizon, Z. Williamson, O. Ciobotaru.
  \emph{PLONK.} ePrint 2019/953.
\bibitem{groth16} J. Groth. \emph{On the size of pairing-based
  non-interactive arguments.} EUROCRYPT 2016.
\bibitem{gkr} S. Goldwasser, Y. T. Kalai, G. N. Rothblum.
  \emph{Delegating computation: interactive proofs for muggles.}
  STOC 2008.
\bibitem{circle} U. Hab\"ock, D. Levit, S. Papini. \emph{Circle
  STARKs.} ePrint 2024/278.
\bibitem{zkcnn} T. Liu, X. Xie, Y. Zhang. \emph{zkCNN.} CCS 2021.
\bibitem{dps} J.-G. Dumas, C. Pernet, A. Sedoglavic.
  arXiv:2506.13242; A. Novikov et al. \emph{AlphaEvolve.}
  arXiv:2506.13131.
\bibitem{ezkl} EZKL. \url{github.com/zkonduit/ezkl}.
\bibitem{plonky2} Polygon Zero. \emph{Plonky2.} Whitepaper, 2022.
\bibitem{risc0} RISC Zero. \emph{zkVM documentation}, 2023.
\bibitem{nova} A. Kothapalli, S. Setty, I. Tzialla. \emph{Nova.}
  CRYPTO 2022.
\bibitem{km} M. Kauers, J. Moosbauer. \emph{Flip graphs for matrix
  multiplication.} arXiv:2212.01175.
\bibitem{poseidon} G. Grassi et al. \emph{Poseidon.} USENIX Security
  2021.
\bibitem{adds} G. Sidebottom. \emph{A 55-addition rank-23 scheme for
  $3\times3$ matrix multiplication.} This repository, 2026. DOI
  10.5281/zenodo.21240904.
\bibitem{companions} G. Sidebottom. \emph{Rigidity of the rank-48
  $4\times4$ scheme} and \emph{The Flower.} Companion notes, this
  repository, 2026.
\end{thebibliography}

\end{document}
